\begin{center}
{\Large \textbf{Sinopsis}}
\end{center}
Progresul recent în prelucrarea limbajului natural a fost accelerat de 
dezvoltarea modelelor lingvistice de mari dimensiuni(LLM). Acest tip de arhitectura 
are ca scop învățarea unor reprezentări ale limbajului care pot captura proprietăți 
sintactice, semantice și să învețe relații între cuvinte îndepărtate într-un mod efficient
și scalabil. În ciuda progresului acestor modele, modul în care ele 
învață și manipulează informația este doar parțial înțeles. \\ \\
%
Această limitare este cu atât mai relevantă în arhitecturile profunde care se bazează pe 
abordări diferite de a înmagazina informația. Modelele de tip transformer procesează
secvențe de text cu ajutorul mecanismului de atenție, care folosește informația din tokeni în fiecare
strat intern, fără o structură definită pentru memorie. Pe de altă parte, xLSTM o arhitectură 
recurentă a fost propusă recent de către Beck et al.\cite{beck2024xlstmextendedlongshortterm} ca o 
alternativă la modelele transformer, extinzând versiunea clasică a celulei LSTM. \\ \\
%
Această lucrare propune o analiză a arhitecturii xLSTM printr-un set de experimente menite să 
ajute la înțelegerea felului în care informația este reprezentată, stocată si transformată în cadrul modelului. 
Mai mult, evaluăm practicalittea acestei arhitecturi într-un scenariu de tip
intrebare-raspuns, o utilizare foarte frecventă a modelelor LLM. \\ \\
%
Se propune observarea evoluției reprezentărilor latente intermediare, precum și comparația cu cele produse de 
un model Transformer pentru a evalua similaritatea dintre ele. 
Mai mult, analizăm comportamentul stării de memorie pe parcursul procesării textului pentru a 
înțelege felul în care tokeni individuali influențează adăugarea 
și eliminarea de informație în memorie. \\ \\ 
%
Concluzionând, evaluăm stocarea în cache a stării de memorie a modelului, exploatând dimensiunea
fixă a acesteia. Prin reutilizarea stării de memorie după procesarea contextului, modelul evită reprocesarea contextului
și îmbunătățește viteza de generare a răspunsurilor direct proporțional cu lungimea contextului. 
Această metodă este evaluată în comparație cu utilizarea în aceleași condiții a unui model Mistal-7B cu scopul de a examina
avantajele și dezavantajele privind eficiența inferenței și performanța modelului în cadrul experimentului.